\documentclass[12pt]{article}
\usepackage[utf8]{inputenc}
\usepackage[T1]{fontenc}
\usepackage{hyperref}
\title{Bora Language Resources for K-12 Curriculum Development}
\author{Compiled by Claude Sonnet 5.5 (claude-sonnet-5-5) for Andr\'as Kornai}
\date{}
\begin{document}
\maketitle

\section{Introduction}
Bora (ISO 639-3 \texttt{boa}, Glottocode bora1263) is a Boran language of Peru and Colombia. Search results give estimates of roughly 2,000--3,000 speakers, most in Peru and about 500 in Putumayo, Colombia; Mira\~na is treated as a closely related dialect. This report lists only what turned up in web searches. Little pedagogical material was found.

\section{Primary Educational Resources}
No primers, textbooks or school materials for Bora were found in the searches. The Thiesen and Weber grammar below may serve as a source for materials development.

\section{Linguistic and Reference Materials}
\begin{itemize}
\item Wesley Thiesen and David Weber, \emph{A Grammar of Bora with Special Attention to Tone} (SIL International, 2012; about 550 pages): \url{https://mitpressbookstore.mit.edu/book/9781556713019}
\item Wesley Thiesen, \emph{Gram\'atica del idioma bora} (Serie Ling\"u\'istica Peruana 38, 2nd ed., 2008), Glottolog record: \url{https://glottolog.org/resource/reference/id/314538}
\item Overview of the language: \url{https://en.wikipedia.org/wiki/Bora_language}
\end{itemize}

\section{Religious and Community Resources}
None found in the searches.

\section{Digital and Audio Resources}
\begin{itemize}
\item Endangered Languages Project entry: \url{https://endangeredlanguages.com/node/97114}
\item Rosetta Project entry, which mentions that a dictionary exists; I did not find the dictionary itself: \url{https://rosettapanglossia.longnow.org/index.php?title=Boa}
\end{itemize}

\section{Curriculum Development Context}
The language is reported as not known to be taught in schools. Any curriculum would have to be built from the grammar and from new work with speakers. The resources found are few, and I verified no audio, dictionary or school programme.
\end{document}